% ============================================================
%  ZoPhysicist content file.
%  To add a new lesson: copy this file, rename it (e.g. waves-intro.tex),
%  fill in the header below, write your LaTeX after the header, commit.
%  The site rebuilds and deploys itself automatically.
% ============================================================
%
%  HEADER  (every line starts with  % !ZP  — keep that exactly)
%
% !ZP title   = Put the lesson title here
% !ZP type    = lecture
% !ZP topic   = Kinematics
% !ZP summary = One short sentence shown on the card.
% !ZP tags    = motion, example
% !ZP order   = 10
%
%  FIELD GUIDE
%   title   (required)  Shown as the page heading and in search.
%   type    (required)  One of:  lecture | definition | problem | solution
%   topic   (required)  ANY text. Typing a new topic name automatically
%                       creates that section in the site navigation.
%   summary (optional)  Short blurb on the browse card.
%   tags    (optional)  Comma-separated, used by search.
%   order   (optional)  Number for sorting within a topic (low = first).
%   pair    (optional)  Link a problem and its solution: give BOTH the
%                       problem file and the solution file the SAME pair
%                       value, e.g.  % !ZP pair = kin-01
%   topicorder (optional) Number controlling where the topic itself sits
%                       in the sidebar (put it on any one file of the topic).
%
%  Mizo version (optional): anywhere in the body, put a line that is exactly
%       %%% MIZO %%%
%  Everything ABOVE it is English; everything BELOW it is the Mizo version.
%  The site's EN/MIZO switch then shows the right one. Omit it for English-only.
%
% ------------------------------------------------------------
%  YOUR LATEX STARTS BELOW. Write normally. A preamble /
%  \documentclass is optional — body-only is fine.
% ------------------------------------------------------------

\section*{Put the lesson title here}

Write your explanation in normal LaTeX. Inline maths like $E = mc^{2}$ and
displayed maths both work:

\begin{equation}
  v^{2} = u^{2} + 2 a s
\end{equation}

\begin{itemize}
  \item Lists work.
  \item \textbf{Bold} and \emph{italic} work.
\end{itemize}

Tables, \texttt{align} environments, blockquotes and so on are all supported.

%%% MIZO %%%

\section*{Mizo title (optional)}

Mizo version goes here. Equations need no translation.
